\documentclass[10pt,a4paper]{amsart}
\usepackage[margin=19mm]{geometry}
\usepackage{amsmath,amssymb,mathtools}
\usepackage[scr=rsfs,cal=euler]{mathalfa}
\usepackage{xfrac}
\usepackage[unicode,hidelinks,bookmarks=false]{hyperref}
\newcommand{\ie}{i.e.\ }
\newcommand{\cf}{cf.\ }
\newcommand{\resp}{respectively\ }
\newcommand{\Subsection}{\subsection}
\providecommand{\nobreakdash}{\nobreak\hbox{-}}
\setlength{\emergencystretch}{2em}
\multlinegap0pt
\usepackage{bm}
\usepackage{bbm}
\usepackage{dsfont}
\usepackage{xspace}
\usepackage{cancel}
\usepackage{mathtools}
\usepackage{amsthm}
\makeatletter
\@ifundefined{Theorem}{\newtheorem{Theorem}{Theorem}}{}
\@ifundefined{Lemma}{\newtheorem{Lemma}[Theorem]{Lemma}}{}
\makeatother
\makeatletter
\newcommand{\Inn}{{\mathrm Inn}}
\@ifundefined{Proof}{\newenvironment{Proof}{{\bf Proof.}}{\hfill $\blacksquare$}}{}
\makeatother
\title[The Cayley graph of a quandle]{The Cayley graph of a quandle}
\author{David Dol\v{z}an, Bogdana Oliynyk}
\date{Source: arXiv:2604.17011v1 (CC BY 4.0)}
\begin{document}
\maketitle

\noindent\emph{Source and licence.} The Cayley graph of a quandle. Version arXiv:2604.17011v1 (CC BY 4.0), distributed under the Creative Commons Attribution 4.0 International licence, as recorded in the arXiv licence metadata for this exact version and shown on the abstract page https://arxiv.org/abs/2604.17011v1. Full rights notice: licenses/arxiv-2604.17011-CC-BY-4.0.txt.\par\smallskip

\noindent\textit{Editorial scope.} Counted unit: the Theorem whose source label is \texttt{orbits} in the article (source0.tex lines 891--905, source label \texttt{orbits}), delivered together with the article's own complete original proof of it (source0.tex lines 907--966, one \texttt{proof} environment of 60 lines). No mathematics is written, repaired or re-derived: statement and proof are the article's own text line for line. Required context retained uncounted: inng (lines 870--876). Every definition, display and label of those lines is retained unchanged. Omitted from this package and not redistributed: 1--869, 877--890, 906--906, 967--1019. Nothing inside a retained excerpt is omitted or altered: every retained body is delivered exactly as the article's lines are written, so no omission receipt of any kind accompanies this package. Citations: the retained excerpts carry no citation command at all, so no bibliography entry of the article is transcribed and no reference is redistributed. Reference adaptation: none was needed and none was made; every reference inside the delivered unit resolves inside it, so no unresolved reference or citation is produced. Printed slips kept literal: no printed word, symbol, hypothesis or number of the counted statement or of its proof was altered. Layout and class: the article's own class is \texttt{amsart} with packages including amsopn, amsthm, amsfonts, amsmath, amssymb, amscd, babel, nicematrix, mathtools, tikz; the wrapper is the lane's pinned \texttt{amsart} editorial wrapper at 10pt on A4, which restates the theorem-like environments the retained text needs and adds the article's own macro definitions that the retained text needs (\texttt{\textbackslash Inn}), with extra wrapper packages (bm, bbm, dsfont, xspace, cancel, mathtools). The delivered numbering is the unit's own. Assets: no retained span contains a figure environment, a graphics-inclusion command, a PDF-page inclusion, a table or a TikZ picture; no image file is copied into this package, the package declares no asset dependency and no image file is redistributed with it. Licence: version arXiv:2604.17011v1 of this article is distributed under the Creative Commons Attribution 4.0 International licence, as recorded in the arXiv licence metadata for this exact version and shown on the abstract page https://arxiv.org/abs/2604.17011v1. A full rights notice is delivered at licenses/arxiv-2604.17011-CC-BY-4.0.txt. Characters: every retained excerpt is pure ASCII, so the delivered engine, XeLaTeX with its default Latin Modern text font, reports no missing character.
\begin{Lemma}
\label{inng}
Let $Q$ be a finite quandle and $x \in Q$. Then
\[
\mathcal O^+(x)=\{g(x) : g \in \Inn(Q)\}.
\]
\end{Lemma}

\begin{Theorem}
\label{orbits}
Let $G$ be a finite group and $h\in G$. Let $\varphi:G\to G$ be the inner automorphism given by
$\varphi(g)=hgh^{-1}$ and consider the generalized Alexander quandle $A_\varphi(G)$.
Let
\[
N=\Big\langle\,[h,g]\ \Big|\ g\in G\,\Big\rangle
\]
be the subgroup generated by the commutators
$[h,g]=hgh^{-1}g^{-1}$.
Then 
\[
\mathcal O^{+}(x)=xN \qquad \text{for all } x\in G.
\]
\end{Theorem}

\begin{proof}
Note that we have $[h,g]^{-1}=[g,h]$ for all $g\in G$.
For $x\in G$ we have
\[
x\,[h,x^{-1}]^{-1} \;=\; x\,(x^{-1}hx h^{-1}) \;=\; hxh^{-1}.
\]
Therefore, for all $x,y\in G$,
\[
x\rhd y
= hxy^{-1}h^{-1}y
= (hxh^{-1})\,(hy^{-1}h^{-1}y)
= x\,[h,x^{-1}]^{-1}\,[h,y^{-1}] \in xN.
\]
Thus every right translation, as well as any composition of right translations, maps $xN$ into itself.  Hence
\[
\mathcal O^{+}(x) \subseteq xN.
\]

Let us now prove that we also have  $xN \subseteq \mathcal O^{+}(x)$ for every $x \in G$.
Denote the identity of $G$ by $e$. Then
\[
R_e(x)=x\rhd e=hxe^{-1}h^{-1}e=hxh^{-1},
\]
so $R_e$ is a permutation of the finite set $G$. Let $m\ge 1$ be the (finite) order of $R_e$ in
$\mathrm{Sym}(G)$, so $R_e^m=\mathrm{id}$ and hence $R_e^{-1}=R_e^{m-1}$.
Fix $u\in G$ and consider the map
\[
T_u := R_u\,R_e^{m-1} \in \Inn(A_\varphi(G)).
\]
For any $x\in G$, we have
\[
T_u(x)=R_u(R_e^{-1}(x)).
\]
But $R_e^{-1}(x)=h^{-1}xh$, and therefore
\[
T_u(x)
= (h^{-1}xh)\rhd u
= h(h^{-1}xh)u^{-1}h^{-1}u
= x\,(hu^{-1}h^{-1}u)
= x\,[h,u^{-1}].
\]
So $T_u$ acts as right multiplication by the commutator $[h,u^{-1}]$.
However, $\{u^{-1} : u\in G\}=G$, so the elements $[h,u^{-1}]$
($u\in G$) generate $N$. Consequently, for every $n\in N$ there exist $u_1,\dots,u_r\in G$ and
$\varepsilon_1,\dots,\varepsilon_r\in\{\pm1\}$ such that
\[
n=[h,u_1^{-1}]^{\varepsilon_1}\cdots [h,u_r^{-1}]^{\varepsilon_r}.
\]
Hence $xn=T_{u_r}^{\varepsilon_r}\ldots T_{u_2}^{\varepsilon_2}T_{u_1}^{\varepsilon_1}(x)$. 
Therefore, the right-multiplication map $\rho_n:x\mapsto xn$ lies in $\Inn(A_{\varphi}(G))$, so
$xN \subseteq \{g(x) : g \in \Inn(A_{\varphi}(G))\}=\mathcal O^+(x)$ by Lemma \ref{inng}.
%\[
%\rho_n \;=\; T_{u_1}^{\varepsilon_1}\cdots T_{u_r}^{\varepsilon_r}.
%\]
%Each $T_{u_i}$ is a permutation of the finite set $G$, hence has finite order $t_i$, so
%$T_{u_i}^{-1}=T_{u_i}^{t_i-1}$ is a \emph{positive} power. Therefore $\rho_n$ can be written as %a composition of positive powers of the maps $T_{u_i}$.
%Applying this composition to $x$ gives $\rho_n(x)=xn\in\mathcal O^{+}(x)$.
%Hence $xN\subseteq \mathcal O^{+}(x)$.
This proves that $\mathcal O^{+}(x)=xN$.
\end{proof}

\end{document}
